% §14.4 False-Positive Rate Analysis
%
% Cross-reference resolution status:
%
% DEFINED here:
%   \label{sec:fpr-analysis}                     -- this section
%   \label{prop:bhf-fpr}                         -- the BHF FPR proposition
%
% ALREADY DEFINED (chapter 5):
%   \Cref{def:bloom-filter}                      -- chapter 5 §5.2
%
% ALREADY DEFINED (chapter 8):
%   \Cref{thm:entropy-floor}                     -- chapter 8 §8.3
%
% ALREADY DEFINED (chapter 14, earlier sections):
%   \Cref{def:bhf}                               -- §14.2
%   \Cref{sec:bhf-construction}                  -- §14.2
%   \Cref{alg:bhf-insert-query}                  -- §14.2 pseudocode
%   \eqref{eq:bhf-r}                             -- §14.2 fingerprint width
%   \eqref{eq:bhf-insert}                        -- §14.2 build pass
%   \Cref{sec:space-analysis}                    -- §14.3
%   \Cref{thm:bhf-optimality}                    -- §14.3
%
% Citations:
%   \cite{carter-wegman-1979}                    -- universal hashing (in bib)
%
% §14.4 completes the FPR argument previewed at the end of §14.2.
% Sets up the question for x not in A; argues the two cases for phi(x)
% (collides with some a in A, or lands outside the image); bounds FPR
% by 2^{-r}; states the proposition; discusses the uniformity
% assumption on g; verifies the bound at r=7.

\section{False-Positive Rate Analysis}
\label{sec:fpr-analysis}

\Cref{sec:bhf-construction} previewed the false-positive analysis at
the end of \Cref{def:bhf}: for $x \notin A$ the BHF query reduces to
comparing two independent uniform $r$-bit strings, and the
comparison succeeds with probability $2^{-r}$. The present section
makes that argument careful. The setup is small; the argument is two
cases; the bound is $\fprate \le 2^{-r}$. The hash-function
assumption that the bound rests on, uniformity of the fingerprint
hash $g$ on $\{0, 1\}^r$, is named and connected to the universal
hashing literature. A worked verification at $r = 7$ closes the
section and ties back to the worked example of
\Cref{sec:bhf-construction}.

\subsection*{Setup}

Fix a designated set $A \subseteq \Universe$ with $|A| = n$, a
perfect hash $\phi : A \to [n]$ extended in some manner to all of
$\Universe$, and a fingerprint hash $g : \Universe \to \{0, 1\}^r$
with $r = \lceil -\log_2 \fprate \rceil$ from \eqref{eq:bhf-r}. The
fingerprint table $F$ has been populated by the build pass
\eqref{eq:bhf-insert}: $F[\phi(a)] = g(a)$ for every $a \in A$. Pick
a query point $x \notin A$. The question the section answers is:
\[
  \Prob{\textsc{Query}(x) = \top \Given x \notin A} \;=\; ?
\]
where the probability is taken over the random choice of the
fingerprint hash $g$ (the perfect hash $\phi$ is fixed by the build
phase and is not randomized at query time). The query algorithm of
\Cref{alg:bhf-insert-query} reads slot $\phi(x)$ of the table,
computes $g(x)$, and returns $\top$ iff $F[\phi(x)] = g(x)$. The
question reduces to: when is $F[\phi(x)] = g(x)$ for a non-member
$x$?

\subsection*{The argument: two cases for \texorpdfstring{$\phi(x)$}{phi(x)}}

The perfect hash $\phi$ was constructed to be injective on $A$, but
it is defined on all of $\Universe$ by some extension rule (a
detail of the perfect-hash construction; for the FPR analysis only
the resulting image matters). When $x \notin A$, two structural
situations arise.

\smallskip
\noindent\emph{Case (a): $\phi(x)$ lands in the image $\phi(A) \subseteq
[n]$.} The perfect hash maps $x$ to some slot $i \in [n]$ that
coincides with $\phi(a^*)$ for the unique $a^* \in A$ with $\phi(a^*)
= i$ (uniqueness by injectivity of $\phi$ on $A$). The slot's
contents are $F[i] = g(a^*)$, set by the build pass. The query
compares this to $g(x)$, and returns $\top$ iff $g(a^*) = g(x)$.
Since $x \neq a^*$ (because $x \notin A$ and $a^* \in A$), and since
$g$ is independent of the choice of $a^*$ versus $x$, the probability
of a fingerprint collision is exactly $\Prob{g(a^*) = g(x)}$. Under
the uniformity assumption on $g$, the two values $g(a^*)$ and $g(x)$
are independent uniform random strings in $\{0, 1\}^r$, and their
collision probability is $2^{-r}$.

\smallskip
\noindent\emph{Case (b): $\phi(x)$ lands outside the image $\phi(A)$.}
The perfect hash's extension may map some $x \notin A$ to slots that
no element of $A$ uses, or the extension may flag such inputs as
out-of-range. The two standard implementations behave the same way
for FPR purposes. If the extension flags the input as out-of-range,
the query returns $\bot$ directly (no false positive). If the
extension maps $x$ to a slot $i \in [n] \setminus \phi(A)$, that
slot was never written by the build pass; the slot's contents are
either an explicit sentinel (designed to mismatch every fingerprint)
or an arbitrary residual value, depending on how the construction
initializes unused slots. In either implementation, the slot's
contents are independent of $g(x)$, so the fingerprint comparison
succeeds with probability at most $2^{-r}$ (and zero in the sentinel
case). Across the two implementations, case (b) contributes at most
$2^{-r}$ to the false-positive rate, and contributes zero in the
sentinel variant the BHF assumes by default.

The two cases together bound the false-positive rate: case (a) is the
dominant contributor and case (b) adds at most an equal amount,
which in the sentinel-initialized variant collapses to zero. The
worst case is when every $x \notin A$ falls into case (a), giving
$\fprate = 2^{-r}$ exactly; the best case is when the perfect-hash
extension flags case (b) as out-of-range, removing the contribution
from non-image queries entirely. Both bracket the central result.

\subsection*{The bound}

Combining the two cases gives the headline inequality.

\begin{proposition}[BHF false-positive rate]
\label{prop:bhf-fpr}
Let the BHF of \Cref{def:bhf} be constructed for a designated set
$A$ with parameters $\phi$, $g$, and $r = \lceil -\log_2 \fprate
\rceil$. Assume the fingerprint hash $g$ is uniformly distributed on
$\{0, 1\}^r$ over the query point. For every $x \notin A$,
\[
  \fprate_{\text{BHF}} \;=\; \Prob{\textsc{Query}(x) = \top \Given
  x \notin A} \;\le\; 2^{-r}.
\]
With $r = \lceil -\log_2 \fprate \rceil$, the bound implies
$\fprate_{\text{BHF}} \le 2^{-r} \le \fprate$. The BHF meets the
designer's target false-positive rate by construction.
\end{proposition}

The proposition is the second half of the optimality story. The
first half, in \Cref{thm:bhf-optimality}, was that the BHF
achieves the entropy floor of \Cref{thm:entropy-floor} up to
integer rounding plus a sublinear correction; the present
proposition is that it does so while meeting the FPR target. The
two together certify the BHF as both space-optimal and accuracy-
correct: it spends the minimum number of bits per element (up to
rounding) and delivers the false-positive rate the designer asked
for.

The inequality in the proposition is tight in the sentinel-free
variant of the construction (when case (b) maps to arbitrary
residual contents, the upper bound $2^{-r}$ is achieved for every
$x \notin A$) and strict in the sentinel variant (when case (b)
contributes zero, the FPR is $2^{-r}$ weighted by the probability
that $\phi(x) \in \phi(A)$, which is strictly less than $1$ when
the universe extends beyond $A$). The sentinel variant is the
canonical implementation and is what most published BHF
constructions use; the analysis proceeds with the upper-bound
inequality regardless, because the bound is what the designer's
target requires.

The chained inequality $2^{-r} \le \fprate$ that the construction
delivers is not an accident of the algorithm but a direct
consequence of the choice $r = \lceil -\log_2 \fprate \rceil$. The
ceiling rounding ensures $r \ge -\log_2 \fprate$, hence $2^{-r}
\le \fprate$. The BHF over-achieves the target by the fractional
part of $-\log_2 \fprate$ (between $0$ and $1$ bit), the same
rounding gap \Cref{sec:space-analysis} measured on the space side.
Accuracy and space pay the rounding cost together: one integer bit
of fingerprint width buys one integer bit of FPR precision, and the
real-valued target falls between two such bits in the generic case.

\subsection*{Independence of the fingerprint hash}

The proposition's hypothesis, that $g$ is uniformly distributed on
$\{0, 1\}^r$ over the query point, deserves comment. In a worst-case
analysis where the adversary chooses $x$ to maximize the
fingerprint collision probability, the hypothesis can fail: an
adversary who knows $g$ can construct $x$ with $g(x) = g(a^*)$ for
some chosen $a^* \in A$, driving the per-query FPR to $1$. The
uniformity assumption protects against the average case, not the
adversarial case. The average-case framing is appropriate for
applications where query points are drawn from a distribution
independent of the implementation's choice of $g$, which covers
most data-structure use cases.

The hash-function literature has formalized the uniformity
requirement under several names. \emph{Pairwise independence} (also
called $2$-independence) of $g$ is sufficient for the proposition:
if $g(a^*)$ and $g(x)$ are pairwise independent for any pair
$(a^*, x)$, the collision probability is $2^{-r}$ by the standard
hash-function argument. \emph{Universal hashing} in the sense of
\textcite{carter-wegman-1979} delivers pairwise independence by
construction: the universal-hash family contains hash functions
$g$ such that for every pair $a^*, x$ with $a^* \neq x$,
$\Prob{g(a^*) = g(x)} \le 1 / |\text{codomain}|$, where the
probability is over the random choice of $g$ from the family. The
BHF construction selects $g$ from a universal family at build time
and treats the selection as fixed for the data structure's
lifetime. The collision probability bound then applies pointwise
over query points, recovering the proposition's $2^{-r}$ bound on
the per-query FPR for any non-adversarial query distribution.

In practice the fingerprint hash $g$ is implemented as a wide
cryptographic hash (SHA-256, BLAKE3) truncated to $r$ bits, or as a
high-quality non-cryptographic hash (MurmurHash3, xxHash) with
similar uniformity properties. The uniformity assumption is
empirically validated by the absence of detectable correlations in
the truncated output across realistic input distributions. The
adversarial case, in which the implementation's hash function is
exposed to an attacker choosing queries to maximize the FPR, is
handled by salting $g$ with a build-time-random secret, denying the
attacker advance knowledge of which hash function is in use.
\Cref{ch:not-bernoulli} returns to the adversarial setting in
broader context; for the present chapter the non-adversarial
uniformity assumption is the operating hypothesis.

\subsection*{Worked verification at \texorpdfstring{$r = 7$}{r = 7}}

The worked example of \Cref{sec:bhf-construction} chose $\fprate =
0.01$ and computed $r = \lceil -\log_2 0.01 \rceil = 7$. The
proposition predicts $\fprate_{\text{BHF}} \le 2^{-7}$. Computing,
\[
  2^{-7} \;=\; \frac{1}{128} \;=\; 0.0078125,
\]
which is strictly smaller than the target $\fprate = 0.01$ by a
factor of $0.0078 / 0.01 \approx 0.78$, or about $22\%$ below the
target. The BHF over-achieves the target by this fractional amount,
the rounding slack identified above. The construction is correct
for the target: every non-member query has at most a $0.78\%$ chance
of returning a false positive, well under the $1\%$ target the
designer set.

The verification can be sharpened slightly by recognizing that the
target $\fprate = 0.01$ is awkward in base $2$: $-\log_2 0.01 =
6.644$, so $\lceil 6.644 \rceil = 7$ and the rounding loses about
$0.36$ bits per element on the space side (the gap measured in
\Cref{sec:space-analysis}) and gains about $0.22 \cdot \fprate$ on
the FPR side. The two are related: the FPR over-achievement is
$1 - 2^{-r} / \fprate = 1 - 2^{6.644 - 7} = 1 - 2^{-0.356} \approx
0.219$, which matches the $22\%$ figure. The rounding behaves
consistently on both sides of the trade-off, as the construction
demands.

If the target were chosen to land cleanly at an integer bit
boundary, the rounding slack would vanish on both sides. For
$\fprate = 2^{-7} = 0.0078125$, the BHF achieves exactly the target
with $r = 7$, no slack on either FPR or space. For $\fprate =
2^{-10} = 0.000976$, the BHF achieves exactly the target with
$r = 10$. The construction is most economical at power-of-two FPR
targets; awkward targets pay the rounding price, which is at most
one bit per element and at most a factor of $2$ in FPR. Either
price is small compared to the Bloom filter's multiplicative $\log_2
e \approx 1.44$ overhead at the same target.

The verification closes the loop opened in \Cref{sec:bhf-construction}.
The construction works: it satisfies the FNR-zero guarantee
unconditionally, and it satisfies the FPR-at-most-$\fprate$ target
under the uniformity assumption on $g$ that universal hashing
provides. The bit budget is $\lceil -\log_2 \fprate \rceil$ plus the
sublinear perfect-hash overhead, matching the entropy floor up to
the integer rounding both the FPR and the space share.
\Cref{sec:construction-time} turns from accuracy and space to the
remaining performance dimension, build time and query latency, and
the static-versus-streaming trade-off that distinguishes the BHF
from the Bloom filter as a deployable data structure.
